% Apartado 4: seguridad (el LLM no autoriza).
\section{Seguridad: el LLM no autoriza}\label{sec:seguridad}

La seguridad es simulada, porque no hay un banco real detrás, pero el diseño es el de un banco. Aunque alguien engañe al
modelo, el modelo no tiene forma de mover dinero ni de ver datos ajenos.

\subsection{Identidad}\label{sec:identidad}

El cliente entra con usuario y contraseña. La contraseña se guarda con hash \textbf{scrypt} y sal
(\ruta{salida/seguridad/hasheador.py}). La sesión va en una cookie \cod{HttpOnly} y \cod{SameSite=Strict}, y el servidor
solo guarda la huella sha256 del token.

Las contraseñas de demostración no están en el código. Se leen de \cod{CONTRASENA\_<USUARIO>} en \cod{.env} y, si no
existe, son el usuario seguido de 2026. El administrador no tiene contraseña por defecto, y sin \ruta{CONTRASENA_ADMIN}
(12 caracteres como mínimo) ese perfil no existe. El administrador solo gestiona la IA y los tokens. No tiene cuenta ni
chat, y los clientes reciben un 403 si piden \cod{/ajustes}.

Tras 5 fallos de un usuario, o 100 fallos en 15 minutos desde un mismo origen, el acceso se bloquea 5 minutos
(\ruta{dominio/acceso.py}). El intento se reserva antes de verificar la contraseña, bajo un cerrojo. Sin esa reserva,
40 intentos en paralelo llegaban a probar 160 contraseñas. El mensaje es el mismo para un usuario que no existe y para
una contraseña errónea.

\subsection{Aislamiento por cliente}\label{sec:aislamiento}

Todo puerto y caso de uso recibe el \ruta{cliente_id} de la sesión, nunca del modelo. Las herramientas reciben el
\cod{Cliente} ya autenticado y su esquema no tiene ningún campo para elegir otro. En SQL, las escrituras llevan
\cod{WHERE id = ? AND cliente\_id = ?}, y el SQL que escribe el modelo corre en la copia en memoria del
apartado~\ref{sec:nlsql}. Una operación de otro cliente responde 404, igual que una que no existe.

\subsection{Operaciones con OTP: el flujo de un Bizum}\label{sec:otp}

\ruta{transferencia_bizum} \textbf{no envía dinero}. Valida el importe, el saldo y los topes, crea una operación pendiente
y manda un código de 6 dígitos a un ``móvil simulado'', que es un canal distinto del chat. El cliente lo teclea en la
tarjeta de confirmación, que llama a un endpoint aparte. El código nunca pasa por el modelo, ni por el historial, ni por
los eventos SSE (figura~\ref{fig:secuencia-bizum}).

\begin{figure}[htbp]
  \centering
  % Diagrama de secuencia del Bizum con OTP, en TikZ: participantes arriba, líneas de vida punteadas y mensajes horizontales.
% Flecha continua: petición. Discontinua: respuesta.
\resizebox{\textwidth}{!}{%
\begin{tikzpicture}[
    font=\sffamily\small,
    participante/.style={draw=unicajaPetroleo, fill=unicajaInfoClaro, rounded corners=2pt,
                         align=center, text width=2.4cm, minimum height=1.1cm, inner sep=3pt},
    actor/.style={participante, fill=unicajaLima!60, draw=unicajaPetroleoOscuro},
    vida/.style={densely dotted, line width=0.6pt, draw=unicajaPetroleo!70},
    peticion/.style={-{Stealth[length=2.4mm]}, line width=0.6pt, draw=unicajaTinta},
    respuesta/.style={-{Stealth[length=2.4mm]}, dashed, line width=0.6pt, draw=unicajaVerde},
    texto/.style={font=\sffamily\footnotesize, align=center, inner sep=1.5pt, fill=white,
                   execute at begin node={\hyphenpenalty=10000\exhyphenpenalty=10000}},
  ]

  % Separación horizontal entre líneas de vida y vertical entre mensajes
  \def\dx{2.9}
  \def\dy{1.2}

  % ---------- Participantes ----------
  \node[actor]        (C) at (0*\dx,0) {Cliente\\[-1pt]{\scriptsize (actor)}};
  \node[participante] (F) at (1*\dx,0) {Navegador};
  \node[participante] (S) at (2*\dx,0) {ServicioChat\\ + Agente};
  \node[participante] (M) at (3*\dx,0) {Modelo (LLM)};
  \node[participante] (B) at (4*\dx,0) {SolicitarBizum};
  \node[participante] (T) at (5*\dx,0) {Móvil simulado};
  \node[participante] (K) at (6*\dx,0) {Confirmar\-Operacion};

  % ---------- Líneas de vida ----------
  \def\vidaFin{-22.5}
  \foreach \p in {C,F,S,M,B,T,K} {
    \draw[vida] (\p.south) -- (\p.south |- 0,\vidaFin);
  }

  % Mensaje entre dos líneas de vida: \mensajeSec{estilo}{origen}{destino}{fila}{texto}
  \newcommand{\mensajeSec}[5]{%
    \draw[#1] (#2.south |- 0,{-#4*\dy-0.4}) -- node[texto, above] {#5} (#3.south |- 0,{-#4*\dy-0.4});
  }
  % Mensaje a sí mismo: \propioSec{participante}{fila}{texto}
  \newcommand{\propioSec}[3]{%
    \draw[peticion] (#1.south |- 0,{-#2*\dy-0.25}) -- ++(0.45,0) -- ++(0,-0.45) -- ++(-0.45,0);
    \node[texto, anchor=west, align=left, text width=2.2cm] at ($(#1.south |- 0,{-#2*\dy-0.47})+(0.55,0)$) {#3};
  }

  % ---------- Mensajes ----------
  \mensajeSec{peticion}{C}{F}{1}{``Envía 5 euros a Dani''\\(voz o texto)}
  \mensajeSec{peticion}{F}{S}{2}{POST /chat\\(cookie de sesión)}
  \mensajeSec{peticion}{S}{M}{3}{historial +\\herramientas}
  \mensajeSec{respuesta}{M}{S}{4}{llamada transferencia\_bizum\\(destinatario, importe)}
  \mensajeSec{peticion}{S}{B}{5}{cliente de la sesión + argumentos}
  \propioSec{B}{6}{busca el contacto, valida importe, saldo y topes}
  \mensajeSec{peticion}{B}{T}{7.8}{SMS con código de 6 dígitos\\(se guarda solo su huella)}
  \mensajeSec{respuesta}{B}{S}{8.8}{operación pendiente (sin código)}
  \mensajeSec{respuesta}{S}{F}{9.8}{SSE operacion + fin\\``He preparado el Bizum…''}
  \mensajeSec{peticion}{F}{T}{10.8}{GET /movil/mensajes}
  \mensajeSec{respuesta}{T}{C}{11.8}{código (en pantalla o leído en voz alta)}
  \mensajeSec{peticion}{C}{F}{12.8}{teclea el código o\\aprueba con huella}
  \mensajeSec{peticion}{F}{K}{13.8}{POST /operaciones/\{id\}/confirmar}
  \propioSec{K}{14.8}{operación propia, pendiente, no caducada, código correcto}
  \mensajeSec{respuesta}{K}{F}{16.8}{comprobante: Bizum enviado, saldo actualizado}

  % ---------- Nota sobre M … K ----------
  \node[draw=unicajaAviso, fill=unicajaAvisoClaro, texto, font=\sffamily\small\bfseries,
        minimum height=0.7cm, inner xsep=6pt]
    at ($(M.south |- 0,-21.9)!0.5!(K.south |- 0,-21.9)$)
    [minimum width=11.2cm]
    {El modelo nunca ve el código ni puede llamar a /confirmar};
\end{tikzpicture}%
}

  \caption{Flujo de un Bizum con OTP. Las flechas continuas son peticiones y las discontinuas, respuestas.}
  \label{fig:secuencia-bizum}
\end{figure}

El código caduca a los 5 minutos, admite 3 intentos y sirve una sola vez. Al tercer fallo la operación se cancela y
responde 423 (\ruta{dominio/operaciones.py}). Las reglas de riesgo están en \ruta{dominio/reglas.py}. Un Bizum no puede
pasar de 1.000~€ por operación ni de 1.500~€ al día, y la tarjeta muestra un aviso si el destinatario es nuevo o el
importe es de 300~€ o más. Confirmar, aprobar y cancelar comparten un cerrojo por cliente. Antes de ponerlo, ocho
aprobaciones en paralelo restaban 80~€ en vez de 10~€.

Con las tarjetas hay una asimetría a propósito. Bloquear no pide código, porque no mueve dinero, va en la dirección
segura y quien ha perdido la cartera tiene prisa. Desbloquear sí lo pide, con el mismo flujo que el Bizum, porque es
justo lo que querría quien la ha robado. El servidor elige la tarjeta (si hay varias, pregunta) y solo bloquea ante una
orden explícita en primera persona. ``A mi madre le han robado la tarjeta'' no bloquea la del cliente.

\subsection{Otras capas}\label{sec:otras-capas}

{\small
\begin{tabularx}{\textwidth}{@{}P{3.6cm}L@{}}
\toprule
\textbf{Capa} & \textbf{Qué hace} \\
\midrule
Un turno a la vez por cliente & Una segunda pregunta mientras hay otra en curso recibe 429. Así no se puede saltar el tope diario con peticiones en paralelo \\
Marca \cod{[auto]} reservada & Los avisos internos del servidor al modelo llevan esa marca. Si un usuario la escribe, en cualquier variante Unicode, el mensaje se rechaza \\
Auditoría & Login, intentos de inyección, OTP, Bizum, bloqueos, consultas SQL y cambios de ajustes (\ruta{salida/sqlite/auditoria.py}). Los números de tarjeta largos se enmascaran \\
Navegador & CSP \cod{script-src \textquotesingle self\textquotesingle} (librerías servidas desde el propio servidor), HSTS, \cod{X-Frame-Options: DENY}, \cod{no-store} y Swagger apagado por defecto \\
Origen de la IP & Se arranca con \cod{-{}-no-proxy-headers} y el código ignora \cod{X-Forwarded-For} salvo con \ruta{PROXY_DE_CONFIANZA=1} \\
Errores & Cada error de dominio tiene su código HTTP y el 500 no filtra detalles internos \\
\bottomrule
\end{tabularx}}

El código pasó tres rondas de revisión de seguridad. Cada fallo que apareció está corregido y tiene su test de
regresión.
